\section{Introduction}
\label{sec:introduction}

Mean-payoff games ask how two players can control the long-run average
reward of an infinite play. In the stochastic version, some transitions
are chosen by chance. This combines strategic choice with the possibility
of moving between regions having different long-run rewards. The input
may describe very small transition probabilities and very large rewards
with only a few binary digits, so an algorithm whose cost depends on
their magnitudes need not be efficient in the input length.

\subsection{The problem and the result}

An instance consists of a finite directed multigraph $G=(V,E)$ with
$V=\{1,\ldots,n\}$, $n\ge1$, and at least one outgoing edge at every
vertex. The vertex set is partitioned as
\[
 V=V_{\max}\sqcup V_{\min}\sqcup V_{\mathrm{ch}}.
\]
Each edge $e$ has an integer reward $w(e)$. At a vertex of $V_{\max}$,
Max chooses the next edge; at a vertex of $V_{\min}$, Min chooses it;
at a vertex $i\in V_{\mathrm{ch}}$, the next edge is drawn from a
prescribed rational probability distribution $p_i$ on the outgoing
edges. Rewards and the numerators and positive denominators of the
probabilities are encoded in binary. The graph and all edge data are
listed explicitly. Loops, parallel edges, zero probabilities, and
unreduced fractions are permitted.

Both players observe the complete finite history, including traversed
edges, and may use arbitrary behavioral strategies with private
randomness. Chance draws have the prescribed conditional law at the
current vertex. A pure positional strategy chooses one outgoing edge at
each vertex of its player, independently of the preceding history.
For a play $\pi=(X_0,e_0,X_1,e_1,\ldots)$ define
\[
 \MP_w(\pi)=\liminf_{T\to\infty}\frac1T\sum_{t=0}^{T-1}w(e_t),
 \qquad
 \val(i)=\sup_\sigma\inf_\tau
           \E_{i,\sigma,\tau}[\MP_w].
\]
Here $\sigma$ and $\tau$ range over Max's and Min's behavioral
strategies, respectively. The expectation is taken after the pathwise
lower limit. If $L$ is the complete binary input length, the decision
problem is to output the set
\[
                         U=\{i\in V:\val(i)\ge0\}.
\]

\begin{theorem}
\label{thm:main}
There are a uniform deterministic Turing machine and an absolute constant
$C$ such that, on every valid input described above, the machine outputs
exactly $U$ in at most
\[
                       2^{C(\log_2(L+2))^2}
\]
bit operations. The bound includes input reading, exact rational
arithmetic, and output. Vertices of value zero are included.
\end{theorem}

\subsection{Historical context}
\label{subsec:history}

Shapley introduced stochastic games with positive stopping probabilities
and characterized their values by a contraction operator
\cite{Shapley1953}. Gillette considered the undiscounted, long-run
average objective \cite{Gillette1957}; Liggett and Lippman supplied
the classical stationary-strategy theorem for perfect-information
time-average games \cite{LiggettLippman1969}. The relation between
discounting and long-run optimization is also classical. Blackwell
showed in finite dynamic programming that one stationary policy is
optimal for every discount factor sufficiently close to one
\cite{Blackwell1962}. Rational discounted values provide the
corresponding stabilization mechanism for perfect-information
stochastic games; see, for example, Gimbert and Zielonka
\cite[Theorem~3]{GimbertZielonka2010}.

The computational problem connects several familiar classes of games.
A simple stochastic game is a finite turn-based reachability game in
which chance chooses fairly between two successors. Condon placed its
threshold problem in $\mathrm{NP}\cap\mathrm{coNP}$
\cite{Condon1992}, and Ludwig gave a randomized subexponential
algorithm \cite{Ludwig1995}. Zwick and Paterson developed
pseudopolynomial algorithms for deterministic mean-payoff games and
a reduction to simple stochastic games \cite{ZwickPaterson1996}.
Andersson and Miltersen later proved polynomial-time Turing
equivalences among computing values and optimal strategies in
simple stochastic, stochastic mean-payoff, and stochastic
discounted-payoff games, with numerical parameters---including the
discount factor---given in binary
\cite[Theorem~1]{AnderssonMiltersen2009}.
Randomized strategy-improvement algorithms have
subexponential expected bounds \cite{Zwick2026}. A different line
obtains pseudopolynomial bounds when the number of random positions
is fixed \cite{BorosEtAl2019,AllamigeonEtAl2025}. Numerical reward
and denominator parameters in such bounds distinguish them from a
bound in the complete binary input length. Deterministic
quasipolynomial algorithms have also been obtained for simple
stochastic games of bounded treewidth \cite{ChatterjeeEtAl2023}.

Theorem~\ref{thm:main} also gives a deterministic quasipolynomial
algorithm for deciding whether the reachability value of a simple
stochastic game is at least $1/2$. Make the target absorbing, with
reward $+1$ on its self-loop, and give every other edge reward $-1$;
any other terminal vertex receives a self-loop of reward $-1$.
For every strategy pair the pathwise mean payoff is
$2\one_{\{\text{target is reached}\}}-1$.
Its expectation is twice the reachability probability minus one,
so taking the strategy extrema proves the reduction. The transformation
increases the input length by at most a constant factor.

An explicit small discount is only one part of the present algorithm.
Andersson and Miltersen already bounded a sufficient discount using
determinant coefficients and polynomial sign stabilization
\cite[Lemma~1]{AnderssonMiltersen2009}; Gaubert, Grand-Cl\'ement,
and Katz give sharper bounds on Blackwell optimality thresholds
\cite{GaubertGrandClementKatz2025}.
Hansen, Miltersen, and Zwick proved that strategy iteration is
strongly polynomial when the discount factor is fixed
\cite{HansenMiltersenZwick2013}. The dependence of their bound on
the inverse distance of that factor from one does not yield a
polynomial bound in its binary encoding. Here the discount has
polynomial bit length, and a recursive comparison algorithm controls
the cost of locating its fixed point.

For games without chance vertices, a deterministic quasipolynomial
algorithm in the complete binary input length is given in
\citet[Theorem~1.1]{OpenAI2026Deterministic}. Its simultaneous labeling
construction is the direct predecessor of the method used here. We
extend that construction to discounted operators, so that chance
transitions enter through the same order and translation properties
as player choices.

\subsection{The proof and its main ideas}

The first part of the proof replaces the mean-payoff problem by one
discounted fixed point. For $0<u<1$, the discounted map takes the maximum,
minimum, or chance average of $u w(e)+(1-u)t_j$ over edges $e:i\to j$.
It is a contraction and has a unique fixed point $v(u)$. An explicit
integer $H$, of polynomial binary length, gives a nonzero-value gap of
at least $1/H$ and a discount $\lambda$ for which
$\|v(\lambda)-\val\|_\infty\le1/(16H)$. Section~\ref{sec:discount}
proves these estimates by elementary bounds on the rational functions
of a fixed positional pair. Its gain and bias then give pathwise
guarantees against every behavioral opponent. This establishes the
payoff convention used here directly.

Finding the discounted fixed point by ordinary contraction iteration
would take a number of steps depending on $1/\lambda$, which can be
exponential in the input length. The second part of the proof treats
any monotone map $F:\R^n\to\R^n$ satisfying
\[
             F(t+c\one)=F(t)+\gamma c\one,
             \qquad 0<\gamma<1.
\]
In an integer box, a subsolution is an integer vector $y$ satisfying
$F_i(y)\ge y_i$ at coordinates above the lower face; a supersolution
satisfies the reverse inequality at coordinates below the upper face.
A subsolution close to an upper face requires a plus label there,
while a supersolution close to a lower face requires a minus label.
The procedure assigns each coordinate a positive mass for each sign.
Its guarantee holds simultaneously for all comparison vectors whose
coordinates away from the opposite face have total mass at most one.

Section~\ref{sec:labeling} adapts the two-pass pivot and mass-reweighting
method of \citet[Sections~4--5]{OpenAI2026Deterministic}.
The crucial translations have a fixed sign: subsolutions move downward
and supersolutions upward. The discounted translation identity preserves
their required inequalities, allowing comparison vectors in a large box
to be used in smaller recursive boxes. Two pivot passes then arrange a
redistribution of the masses that keeps the required comparisons valid.
At each recursive call, at most one child leaves both mass vectors
unchanged. Every other child increases each product of a coordinate's
two masses by a fixed factor. Starting from uniform masses, only
$O(\log(n+1))$ such choices are possible along a path, even when the
coordinate range is exponentially large.

The final part explains why these boundary labels suffice to find the
fixed point. After scaling, we place an integer subsolution $Y$ and an
integer supersolution $Z$ on opposite sides of it, with $Z-Y$ uniformly
small. Neither vector is computed. A label excluding one comparison
vector from a boundary strip permits that boundary to move inward while
retaining both. Section~\ref{sec:application} shrinks the enclosing box
in this way, and the value gap converts the final approximation into
the exact nonnegative-value set.

The reusable part of the argument is the labeling procedure: its
correctness uses monotonicity and the scalar translation identity,
and does not inspect the graph or the transition probabilities. The
two-sided integer comparison vectors provide the link from that
procedure to a discounted game. Quantitative discounting supplies the
required accuracy, while the recursion determines the deterministic
quasipolynomial running time.

Throughout, vector inequalities, extrema, and integer rounding are
coordinatewise, and $\one$ denotes the all-ones vector. The proof uses
only exact arithmetic.
